% Part: lambda-calculus
% Chapter: introduction
% Section: overview

\documentclass[../../../include/open-logic-section]{subfiles}

\begin{document}

\olfileid{lam}{int}{ovr}
\olsection{مرور کلی}

حساب لامبدا را آلونزو چرچ در اوایل دههٔ ۱۹۳۰ در اصل به‌منزلهٔ بنیادی
برای منطق ساخت‌گرا طراحی کرد، و ‌\emph{نه} به‌عنوان مدلی برای توابع
محاسبه‌پذیر. اما خیلی زود نشان داده شد که این حساب با تعریف‌های دیگر
محاسبه‌پذیری، مانند توابع محاسبه‌پذیر تورینگ و توابع بازگشتی جزئی، هم‌ارز
است. اینکه این نتیجه در آغاز اندکی شگفت‌آور بود، این مشخصه‌سازی را
جالب‌تر می‌کند.

نمادگذاری لامبدا راهی مناسب برای اشارهٔ مستقیم به یک تابع با عبارتی
نمادین است که آن را تعریف می‌کند، بی‌آنکه نامی برایش تعریف کنیم. به‌جای
اینکه بگوییم «بگذارید~$f$ تابعی باشد که با~$f(x) = x + 3$ تعریف شده
است»، می‌توان گفت «بگذارید~$f$ تابع~$\lambd[x][(x + 3)]$ باشد». به بیان
دیگر، $\lambd[x][(x+3)]$ صرفاً \emph{نامی} برای تابعی است که سه واحد به
آرگومان خود می‌افزاید. در این عبارت، $x$ متغیری صوری یا جای‌نگهدار است:
همان تابع را می‌توان به همان خوبی با~$\lambd[y][(y + 3)]$ نشان داد. این
نمادگذاری حتی در حضور پارامترهای دیگر نیز کار می‌کند. برای نمونه، فرض
کنید~$g(x, y)$ تابعی دو‌متغیره و~$k$ عددی طبیعی باشد. آنگاه
$\lambd[x][g(x,k)]$ تابعی است که هر~$x$ را به~$g(x, k)$ می‌نگارد.

این شیوهٔ تعریف تابع با یک عبارت نمادین، \emph{انتزاع لامبدایی} نامیده
می‌شود. روی دیگر انتزاع لامبدایی \emph{کاربرد} است: با فرض داشتن تابعی
چون~$f$ (مثلاً تابعی تعریف‌شده بر اعداد طبیعی)، می‌توان آن را بر هر
ارزشی، مانند ۲، \emph{اعمال کرد}. در نمادگذاری متعارف، البته نتیجه را
به‌صورت~$f(2)$ می‌نویسیم.

وقتی انتزاع لامبدایی را با کاربرد ترکیب کنیم چه رخ می‌دهد؟ عبارت حاصل
را می‌توان با «جای‌گذاری» آرگومانِ کاربرد به‌جای متغیر انتزاع‌شده ساده
کرد. برای نمونه،
\[
(\lambd[x][(x + 3)])(2)
\]
را می‌توان به~$2 + 3$ ساده کرد.

تا اینجا کاری جز معرفی نمادگذاری‌های تازه برای مفاهیم متعارف انجام
نداده‌ایم. بااین‌حال، حساب لامبدا گسستی بنیادی‌تر از دیدگاه
مجموعه‌شناختی را نشان می‌دهد. در این چارچوب:
\begin{enumerate}
\item هر چیز بر تابعی دلالت می‌کند.
\item توابع را می‌توان با انتزاع لامبدایی تعریف کرد.
\item هر چیز را می‌توان بر هر چیز دیگری اعمال کرد.
\end{enumerate}
برای نمونه، اگر~$F$ ترمی در حساب لامبدا باشد، همواره فرض می‌شود که
$F(F)$ معنادار است. این چارچوب آزادانه \emph{حساب لامبدای بی‌نوع}
نامیده می‌شود؛ «بی‌نوع» یعنی «هیچ محدودیتی بر اینکه چه چیزی را می‌توان
بر چه چیزی اعمال کرد وجود ندارد».

\begin{digress}
\emph{حساب لامبدای نوع‌دار} نیز وجود دارد که گونه‌ای مهم از نسخۀ بی‌نوع
است. گرچه حساب لامبدای نوع‌دار از جهات بسیاری به نسخۀ بی‌نوع شباهت دارد،
سازگارکردن آن با چارچوب کلاسیک مجموعه‌شناختی بسیار آسان‌تر است و برخی
ویژگی‌هایش کاملاً متفاوت‌اند.

پژوهش دربارهٔ حساب لامبدا در علوم رایانۀ نظری و طراحی زبان‌های
برنامه‌نویسی نقشی محوری یافته است. لیسپ، که جان مک‌کارتی در دههٔ ۱۹۵۰
طراحی کرد، نمونه‌ای آغازین از زبانی است که از این اندیشه‌ها اثر پذیرفت.
\end{digress}

\end{document}
